\textbf{Phần 1. Trắc nghiệm lựa chọn}
\setcounterpageref{ex}{0}

\begin{ex} %[2D2V07]
Tính giá trị của biểu thức: $9 + 4 + 1 + 6$.
\choice
{$19$}
{$18$}
{\True $20$}
{$21$}
\loigiai{
Gộp các số tạo 10: $(9 + 1) + (4 + 6) = 10 + 10 = 20$.
\textbf{Đáp án của câu hỏi là C}
}
\end{ex}

\begin{ex} %[2D2V07]
Trong hộp có $8$ viên bi xanh và $7$ viên bi đỏ. Bạn Nam lấy thêm vào hộp $3$ viên bi vàng. Hỏi trong hộp có tất cả bao nhiêu viên bi?
\choice
{\True $18$ viên bi}
{$17$ viên bi}
{$19$ viên bi}
{$15$ viên bi}
\loigiai{
Tổng số viên bi là: $8 + 7 + 3 = 8 + (7 + 3) = 8 + 10 = 18$ (viên bi).
\textbf{Đáp án của câu hỏi là A}
}
\end{ex}

\begin{ex} %[2D2V07]
Tìm số tự nhiên $x$ biết: $x - 7 = 9$.
\choice
{$15$}
{$17$}
{$18$}
{\True $16$}
\loigiai{
Ta có: $x = 9 + 7 = 16$.
\textbf{Đáp án của câu hỏi là D}
}
\end{ex}

\begin{ex} %[2D2V07]
Mẹ mua về $9$ quả cam, sau đó mẹ mua thêm $8$ quả cam nữa. Mẹ biếu bà $6$ quả cam. Hỏi mẹ còn lại bao nhiêu quả cam?
\choice
{$12$ quả cam}
{\True $11$ quả cam}
{$10$ quả cam}
{$13$ quả cam}
\loigiai{
Tổng số cam mẹ mua: $9 + 8 = 17$ (quả).\nSố cam mẹ còn lại: $17 - 6 = 11$ (quả).
\textbf{Đáp án của câu hỏi là B}
}
\end{ex}

\begin{ex} %[2D2V07]
Điền số thích hợp vào dấu hỏi chấm: $8 + ? = 15$.
\choice
{$6$}
{$8$}
{\True $7$}
{$9$}
\loigiai{
Ta có: $? = 15 - 8 = 7$.
\textbf{Đáp án của câu hỏi là C}
}
\end{ex}

\begin{ex} %[2D2V07]
Lan hái được $9$ bông hoa, Mai hái được nhiều hơn Lan $5$ bông hoa. Hỏi cả hai bạn hái được bao nhiêu bông hoa?
\choice
{\True $23$ bông hoa}
{$22$ bông hoa}
{$14$ bông hoa}
{$24$ bông hoa}
\loigiai{
Mai hái được: $9 + 5 = 14$ (bông hoa).\nCả hai bạn hái được: $9 + 14 = 23$ (bông hoa).
\textbf{Đáp án của câu hỏi là A}
}
\end{ex}

\begin{ex} %[2D2V07]
Cho phép tính: $8 + 6 \dots 9 + 4$. Dấu thích hợp điền vào chỗ chấm là
\choice
{$<$}
{$=$}
{Không so sánh được}
{\True $>$}
\loigiai{
Vế trái: $8 + 6 = 14$. Vế phải: $9 + 4 = 13$. Vì $14 > 13$ nên điền dấu $>$. 
\textbf{Đáp án của câu hỏi là D}
}
\end{ex}

\begin{ex} %[2D2V07]
Một xe khách có $9$ hành khách. Đến bến có $7$ người lên xe và $4$ người xuống xe. Hỏi trên xe lúc này có bao nhiêu hành khách?
\choice
{$11$ người}
{\True $12$ người}
{$13$ người}
{$10$ người}
\loigiai{
Số người trên xe lúc này là: $9 + 7 - 4 = 16 - 4 = 12$ (người).
\textbf{Đáp án của câu hỏi là B}
}
\end{ex}

\begin{ex} %[2D2V07]
Tổng của số lớn nhất có 1 chữ số và số chẵn lớn nhất có 1 chữ số là
\choice
{$18$}
{$16$}
{\True $17$}
{$19$}
\loigiai{
Số lớn nhất có 1 chữ số là 9. Số chẵn lớn nhất có 1 chữ số là 8.\nTổng là: $9 + 8 = 17$.
\textbf{Đáp án của câu hỏi là C}
}
\end{ex}

\begin{ex} %[2D2V07]
Tìm một số biết rằng lấy số đó trừ đi $8$ thì được kết quả bằng $7 + 2$.
\choice
{\True $17$}
{$16$}
{$18$}
{$15$}
\loigiai{
Vế phải: $7 + 2 = 9$. Số đó là: $9 + 8 = 17$.
\textbf{Đáp án của câu hỏi là A}
}
\end{ex}

\begin{ex} %[2D2V07]
Có ba túi kẹo: túi thứ nhất có $6$ cái, túi thứ hai có $7$ cái, túi thứ ba có $5$ cái. Hỏi cả ba túi có bao nhiêu cái kẹo?
\choice
{$17$ cái kẹo}
{$19$ cái kẹo}
{$16$ cái kẹo}
{\True $18$ cái kẹo}
\loigiai{
Tổng số kẹo là: $6 + 7 + 5 = 13 + 5 = 18$ (cái kẹo).
\textbf{Đáp án của câu hỏi là D}
}
\end{ex}

\begin{ex} %[2D2V07]
Điền số thích hợp: $7 + 8 = 10 + \dots$
\choice
{$4$}
{\True $5$}
{$6$}
{$7$}
\loigiai{
Vế trái: $7 + 8 = 15$. Ta có: $10 + 5 = 15$. Vậy điền số 5.
\textbf{Đáp án của câu hỏi là B}
}
\end{ex}

\textbf{Phần 2. Trắc nghiệm đúng sai}
\setcounterpageref{ex}{0}

\begin{ex} %[2D2V07]
Xét tính đúng sai của các phép tính cộng qua 10 trong phạm vi 20:
\choiceTF
{\True Phép tính $9 + 6$ có kết quả bằng $15$.}
{Khi cộng 8 với 5 ta tách 5 thành 3 và 2 để $8 + 3 = 10$.}
{\True Tổng của $7 + 8$ và $8 + 7$ luôn bằng nhau.}
{Tổng của hai số có một chữ số không thể vượt quá 17.}
\loigiai{
\begin{itemize}
    \item Ý a) Đúng: $9 + 6 = 15$. Đúng.
    \item Ý b) Sai: Phải tách 5 thành 2 và 3 để $8 + 2 = 10$, sau đó $10 + 3 = 13$.
    \item Ý c) Đúng: Tính chất giao hoán của phép cộng: $7 + 8 = 8 + 7 = 15$.
    \item Ý d) Sai: Tổng lớn nhất là $9 + 9 = 18 > 17$.
\end{itemize}
\begin{itemize}
    \item a) \textbf{Đúng}.
    \item b) \textbf{Sai}.
    \item c) \textbf{Đúng}.
    \item d) \textbf{Sai}.
\end{itemize}
}
\end{ex}

\begin{ex} %[2D2V07]
Hai bạn An và Bình cùng chơi ném bóng vào rổ. An ném trúng 8 quả, Bình ném trúng 9 quả:
\choiceTF
{\True Cả hai bạn ném trúng tất cả 17 quả bóng.}
{Bình ném trúng ít hơn An 1 quả bóng.}
{\True Nếu An ném trúng thêm 3 quả nữa thì An sẽ nhiều hơn Bình 2 quả.}
{Để cả hai bạn có số bóng bằng nhau thì Bình phải cho An 2 quả.}
\loigiai{
\begin{itemize}
    \item Ý a) Đúng: $8 + 9 = 17$ quả.
    \item Ý b) Sai: Bình ném được 9 quả, An 8 quả nên Bình nhiều hơn An 1 quả.
    \item Ý c) Đúng: Lúc đó An có: $8 + 3 = 11$ quả. $11 - 9 = 2$ quả. Đúng.
    \item Ý d) Sai: Bình chỉ hơn An 1 quả nên không thể cho 2 quả để bằng nhau.
\end{itemize}
\begin{itemize}
    \item a) \textbf{Đúng}.
    \item b) \textbf{Sai}.
    \item c) \textbf{Đúng}.
    \item d) \textbf{Sai}.
\end{itemize}
}
\end{ex}

\textbf{Phần 3. Trả lời ngắn}
\setcounterpageref{ex}{0}

\begin{ex} %[2D2V07]
Tính giá trị của biểu thức: $7 + 6 + 3 + 4$.
\shortans{20}
\loigiai{
\textbf{Đáp số:} $20$\\
$(7 + 3) + (6 + 4) = 10 + 10 = 20$.
}
\end{ex}

\begin{ex} %[2D2V07]
Bác nông dân nuôi 8 con vịt, sau đó bác mua thêm 9 con vịt nữa. Hôm sau bác bán đi 5 con vịt. Hỏi bác còn lại bao nhiêu con vịt?
\shortans{12}
\loigiai{
\textbf{Đáp số:} $12$\\
Tổng số vịt: $8 + 9 = 17$. Còn lại: $17 - 5 = 12$ con.
}
\end{ex}

\begin{ex} %[2D2V07]
Tìm số tự nhiên x biết: $x - 8 = 6 + 3$.
\shortans{17}
\loigiai{
\textbf{Đáp số:} $17$\\
$x - 8 = 9 \Rightarrow x = 9 + 8 = 17$.
}
\end{ex}

\begin{ex} %[2D2V07]
Một cuộn ruy băng dài 19 cm. Bạn Lan cắt lấy một đoạn dài 8 cm, sau đó cắt tiếp một đoạn dài 6 cm. Hỏi cuộn ruy băng còn lại dài bao nhiêu xăng-ti-mét?
\shortans{5}
\loigiai{
\textbf{Đáp số:} $5$\\
Độ dài còn lại: $19 - 8 - 6 = 5\text{ cm}$.
}
\end{ex}

\begin{ex} %[2D2V07]
Tìm một số biết rằng lấy số đó cộng với 7 thì được kết quả bằng 16.
\shortans{9}
\loigiai{
\textbf{Đáp số:} $9$\\
Số đó là: $16 - 7 = 9$.
}
\end{ex}

\begin{ex} %[2D2V07]
Có bao nhiêu cặp số tự nhiên có một chữ số mà tổng của chúng bằng 14?
\shortans{5}
\loigiai{
\textbf{Đáp số:} $5$\\
Các cặp là: $(5, 9), (6, 8), (7, 7), (8, 6), (9, 5)$. Có 5 cặp.
}
\end{ex}
